\chapter{Hintergrund}
    - Technologien beschreiben, die nicht von dir sind, die man aber zum Verständnis der Arbeit braucht: "Was ist das eigentlich" --> so erklären, dass jmd ohne Spezialisierungskenntnis, aber mit Bachelor Informatik es verstehen kann

\section{FPGA}
    Field Programmable Gate Array
    Konfigurierbare Hardwareschaltung --> "Hardwareentwicklung" kann sich verhalten wie ein Addierer, eine ALU, wie eine CPU, ein Timer, oder ein bestimmter Microcontroller
    Besteht aus Logikblöcken, die konfiguriert werden können. Die Logikblöcke wiederum enthalten Input/Output Elemente (zur Kommunikation, bspw. LED oder UART), Switching/Routing-Blöcke zur Verbindung der Bauteile sowie Konfiguraitons-Logikblöcke
    Konfiguration-Logikblöcke generieren aus einem Inputsignal zusammen mit einem Clock- und Resetsignal einen Output durch die Verwendung von Lookup-Tables, Flipflops und Multiplexern. Durch diese Bestandteile wird dann eine Schaltung umgesetzt.
    So kann z.B. jede Schaltung, die nicht zeitabhängig ist/die eine logische Schaltung bzw. deterministisch ist, kann in einer Lookup-Table abgebildet werden. Für zeitabhängige Schaltungen benötigt man Flipflops und Speicherbausteine
    Die Konfiguration der Schaltung wurde im Rahmen der Arbeit immer im SRAM gespeichert, war somit flüchtig und musste bei jedem Neustart des Boards erneut geladen werden

    FPGAs werden mehrschrittig konfiguriert:
        RTL-Code (Verilog/VHDL) --> Analyse/Simulation --> Synthese --> Netzliste wird erstellt (Welche Bauteile bräuchte man theoretisch, analog zu Zwischencode beim Kompilieren von Code)
        --> Place \& Route/Implementierung (welche physischen Bauteile werden verwendet und wie werden diese verbunden?)
        --> Bitstream --> Übertragen/konfigurieren/programmieren des Bitstreams auf FPGA

    Einordnung:
    Flexibler als Microcontroller, FPGA-Schaltung ist modifizierbar, kann das werden, was man möchte. Dadurch ist es schneller und energieeffizienter als Microcontroller, aber auch komplexer als Microcontroller
    kann "echt parallel" arbeiten

    Unterscheiden sich hinsichtlich der Leistung/Mächtigkeit vor allem auch in folgenden Kategorien:
        Anzahl I/O Pins, Anzahl Logikblöcke (Wie viele Lookup-Tables), maximale Taktfrequenz, RAM-Größe, Stromverbrauch

\section{Basys-3 Board}
    FPGA-Board mit Artix-7 FPGA + Peripherie:
        FPGA selber, Buttons, Schalter, 7-Segment-Anzeige, LEDs, VGA-Ausgang, USB-Port, Anschlüsse für Erweiterungen. Für die Arbeit wurden lediglich der FPGA selber und eine LED verwendet
    Basys-3 Board Spezifikationen: 100MHz Clock, 33k logic cells --> 312k LUTs, 1800Kbit RAM
\section{Verilog/Vivado}
    Verilog beschreibt keinen Algorithmus, sondern eine Schaltung. Am weitesten verbreitete RTL-Sprache (RTL)
    Vivado ist eine Entwicklungsumgebung für Hardware-Design
        --> Es enthält Funktionalitäten zum Schreiben und Organisieren von RTL-Code, Synthese- und Implementierungsschritte, Generierung des Bitstreams sowie Möglichkeiten zur zykelakkuraten Simulation der Schaltung

\section{RISC-V/SERV}
    RISC-V ist eine open-source Instruction Set Architecture (ISA), die die "Sprache" des Prozessors festlegt. RISC-V CPUs sind solche, die diese Prozessorinstruktionen verarbeiten können.
    SERV ist eine sehr kleine open-source bit-serielle (1 bit pro Takt) RISC-V CPU.
    Servant ist ein System on Chip (SoC), welches die SERV-CPU mit minimaler Peripherie (Timer, Speicher, 1-bit GPIO-Pin) erweitert, um das Zephyr-Betriebssystem auf der CPU laufen lassen zu können.
    Servant ist am Ende auch das SoC, welches auf dem Artix-7 FPGA synthetisiert wurde.

\section{Verilator}
    Simulator für Verilog-Code: Übersetzt den RTL-Code in ein eigenes, schnelles Verilog-Modell, das dann in C++-Code gewrappt wird und simuliert somit die
    Erreicht durch diesen Zwischenschritt eine sehr schnelle Simulation (und ist mindestens gleich schnell, oft schneller als andere Simulatoren)
    Unterschied zu anderen Simulationen: low-level, gleichzeitig sehr schnell
    Verilator wird wiederum mit fusesoc genutzt, um serv-RTL-Code zu simulieren und ist das zugrundeliegende Tool für die im Rahmen der Bachelorarbeit durchgeführten Simulationen

\section{Zephyr}
    Open-source Real-Time Operating System (RTOS) für FPGAs, sehr basic, sodass es auch auf kleine RISC-V Prozessoren läuft
    Quasi Softwareplattform für FPGAs/Microcontroller
    Alternative zu Zephyr: Buildtooling von SERV --> "bare metal"

\section{fusesoc}
    Open-source Buildtool für Hardwareprojekte, organisiert bspw. Abhängigkeiten und kann bspw. Simulation eines SERV-Kerns umsetzen und Zephyr einbinden




% Überlegen ob noch andere Begrifflichkeiten/Technologien erläutert werden müssten